\DocumentMetadata{
  pdfversion=2.0,
  pdfstandard=ua-2,
  lang=en-US,
  tagging=on,
  tagging-setup={math/setup=mathml-SE,tabsorder=structure}
}
\documentclass[11pt,letterpaper]{article}
\usepackage[margin=0.68in,includefoot]{geometry}
\usepackage{newcomputermodern}
\usepackage{amsmath,amscd,mathtools}
\usepackage{tikz,adjustbox}
\providecommand{\square}{\mdlgwhtsquare}
\usepackage[hidelinks]{hyperref}
\hypersetup{
  pdftitle={Algebra First-Year Exam, January 2026, Part 2},
  pdfauthor={University of Florida Department of Mathematics},
  pdfsubject={Graduate mathematics examination},
  pdfdisplaydoctitle=true,
  bookmarksopen=true
}
\setcounter{secnumdepth}{0}
\setlength{\parindent}{0pt}
\setlength{\parskip}{0.28em}
\setlength{\emergencystretch}{3em}
\setlength{\leftmargini}{2.15em}
\setlength{\leftmarginii}{2.65em}
\setlength{\leftmarginiii}{3em}
\renewcommand{\labelenumi}{\textbf{\theenumi.}}
\pagestyle{empty}
\raggedbottom
\allowdisplaybreaks
\newenvironment{examproblems}[1]{%
  \begin{enumerate}%
  \setcounter{enumi}{#1}%
  \setlength{\topsep}{0.35em}%
  \setlength{\partopsep}{0pt}%
  \setlength{\itemsep}{0.48em}%
  \setlength{\parsep}{0.18em}%
}{\end{enumerate}}
\newenvironment{examsubparts}{%
  \begin{enumerate}%
  \setlength{\topsep}{0.18em}%
  \setlength{\partopsep}{0pt}%
  \setlength{\itemsep}{0.16em}%
  \setlength{\parsep}{0.1em}%
}{\end{enumerate}}
\newenvironment{examinstructions}{%
  \par\begingroup\itshape
}{%
  \par\endgroup\vspace{0.18em}
}
\makeatletter
\renewcommand\section{\@startsection{section}{1}{0pt}%
  {-1.25ex plus -.25ex minus -.1ex}%
  {0.55ex plus .1ex}%
  {\normalfont\large\bfseries}}
\renewcommand{\@maketitle}{%
  \newpage\null\vskip -1em%
  {\footnotesize\bfseries
    \href{https://www.ufl.edu/}{University of Florida}, \href{https://math.ufl.edu/}{Department of Mathematics}\par}%
  \vskip -0.5em%
  \tagstructbegin{tag=Title}%
    {\LARGE\bfseries\href{https://gma.math.ufl.edu/exams/algebra-first-year/}{Algebra First-Year Exam}, January 2026, Part 2\par}%
  \tagstructend%
  \vskip 0.25em%
  \tagmcbegin{artifact}%
    \hrule height 1pt%
  \tagmcend%
  \vskip 1em%
}
\makeatother
\title{Algebra First-Year Exam, January 2026, Part 2}
\author{}
\date{}
\begin{document}
\maketitle
\thispagestyle{empty}
\begin{examinstructions}
Answer 4 questions. Write your solutions clearly in complete English sentences. You may quote theorems (within reason) as long as you state them clearly.
\end{examinstructions}
\begin{examproblems}{0}
\item
Prove that the following polynomials are irreducible in \(\mathbb{Q}[x]\), stating clearly any theorems you wish to apply.
\begin{examsubparts}
\item[\textnormal{(a)}]
\(x^4-2x^3+6x^2-x+5\)
\item[\textnormal{(b)}]
\(x^5-4x+2\)
\item[\textnormal{(c)}]
\(x^4+x^3+x^2+x+1\)
\end{examsubparts}
\item
Let~\(R\) be a commutative ring with \(1\ne 0\).
\begin{examsubparts}
\item[\textnormal{(a)}]
Define what it means for an element \(r\in R\) to be nilpotent.
\item[\textnormal{(b)}]
Prove that the set of all nilpotent elements of~\(R\) forms an ideal (called the nilradical of~\(R\)).
\item[\textnormal{(c)}]
Show that the nilradical is contained in every prime ideal of~\(R\).
\end{examsubparts}
\item
Let \(R=\mathbb{Z}[\sqrt{-3}]=\{a+b\sqrt{-3}\mid a,b\in\mathbb{Z}\}\).
\begin{examsubparts}
\item[\textnormal{(a)}]
Show that the function \(N(a+b\sqrt{-3})=a^2+3b^2\) is multiplicative, that is, \(N(rs)=N(r)N(s)\) for all \(r,s\in R\).
\item[\textnormal{(b)}]
Prove that~\(R\) is not a unique factorization domain.
\item[\textnormal{(c)}]
Is~\(R\) a principal ideal domain? Justify your answer.
\end{examsubparts}
\item
Let~\(R\) be a commutative ring with 1 and let~\(M\) be a module over~\(R\).
\begin{examsubparts}
\item[\textnormal{(a)}]
State what it means for~\(M\) to be a Noetherian~\(R\)-module, and for~\(R\) to be a Noetherian ring.
\item[\textnormal{(b)}]
Prove that if~\(M\) is a Noetherian~\(R\)-module, then every nonempty set of submodules has a maximal element.
\item[\textnormal{(c)}]
Give an example of a non-Noetherian commutative ring with 1.
\end{examsubparts}
\item
Using the rational canonical form, find one representative of each conjugacy class of elements of order 3 in \(\mathrm{GL}(4,\mathbb{Q})\).
\item
Let \(K\subseteq E\) be fields.
\begin{examsubparts}
\item[\textnormal{(a)}]
Define what it means for an element of~\(E\) to be algebraic over~\(K\).
\item[\textnormal{(b)}]
Prove that if \(\alpha,\beta\in E\) are algebraic over~\(K\), then \(K(\alpha,\beta)\) is finite-dimensional over~\(K\).
\item[\textnormal{(c)}]
Deduce that the elements of~\(E\) that are algebraic over~\(K\) form a field.
\end{examsubparts}
\end{examproblems}
\end{document}
